% GENERATED FILE. Edit canonical lessons, then run npm run generate:books.
% canonical-source-hash: fnv1a64:4a14ecf7c3c2c329
% canonical-lessons: KA-C306-moci, KA-C306-kumbara, KA-C306-mestri, KA-C306-plambar, KA-C306-totagara

\chapter{Cobbler, Potter, Mason, Plumber, Gardener}
\label{ch:ka-cobbler-potter-mason-plumber-gardener}
\hlchaptermodality{306}

\begin{chapteropening}
\textbf{By the end of this chapter:} \emph{I can say a cobbler, a potter, a mason, a plumber and a gardener.}

Complete the last lesson of chapter 306: I can say a cobbler, a potter, a mason, a plumber and a gardener.
\end{chapteropening}

\section[\texorpdfstring{\kn{ಮೋಚಿ} (mōci)}{ಮೋಚಿ (mōci)}]{\kn{ಮೋಚಿ} (mōci) — a cobbler}

\label{lesson:KA-C306-moci}



\begin{quote}
Before the new one: say the Kannada for an actor, then the Kannada for a journalist.
\end{quote}

\subsection*{You'll want to know: \kn{ಮೋಚಿ}}
\textbf{\kn{ಮೋಚಿ}} — \emph{mōci} — \textquotedblleft{}a cobbler\textquotedblright{}.

\begin{grammarlens}[title={What you've built}]
One more word for this chapter.
\end{grammarlens}

\subsection*{Your turn}
\begin{itemize}
\raggedright
  \item \emph{Say it:} \emph{mōci}
  \item \emph{Say it:} \emph{mōci}, once more
  \item \emph{Say it:} say \emph{mōci}
  \item \emph{Recall:} say the Kannada for an actor, then the Kannada for a journalist, then say \emph{mōci} again
\end{itemize}

\begin{tcolorbox}[breakable,colback=teal!4,colframe=teal!35!black,title={Before you move on}]
What does \kn{ಮೋಚಿ} mean? (\textquotedblleft{}A cobbler\textquotedblright{}.) Say it once more.
\end{tcolorbox}
\section[\texorpdfstring{\kn{ಕುಂಬಾರ} (kumbāra)}{ಕುಂಬಾರ (kumbāra)}]{\kn{ಕುಂಬಾರ} (kumbāra) — a potter}

\label{lesson:KA-C306-kumbara}



\begin{quote}
Before the new one: say the Kannada for a journalist, then the Kannada for a cobbler.
\end{quote}

\subsection*{You'll want to know: \kn{ಕುಂಬಾರ}}
\textbf{\kn{ಕುಂಬಾರ}} — \emph{kumbāra} — \textquotedblleft{}a potter\textquotedblright{}.

\begin{grammarlens}[title={What you've built}]
One more word for this chapter.
\end{grammarlens}

\subsection*{Your turn}
\begin{itemize}
\raggedright
  \item \emph{Say it:} \emph{kumbāra}
  \item \emph{Say it:} \emph{kumbāra}, once more
  \item \emph{Say it:} say \emph{kumbāra}
  \item \emph{Recall:} say the Kannada for a journalist, then the Kannada for a cobbler, then say \emph{kumbāra} again
\end{itemize}

\begin{tcolorbox}[breakable,colback=teal!4,colframe=teal!35!black,title={Before you move on}]
What does \kn{ಕುಂಬಾರ} mean? (\textquotedblleft{}A potter\textquotedblright{}.) Say it once more.
\end{tcolorbox}
\section[\texorpdfstring{\kn{ಮೇಸ್ತ್ರಿ} (mēstri)}{ಮೇಸ್ತ್ರಿ (mēstri)}]{\kn{ಮೇಸ್ತ್ರಿ} (mēstri) — a mason, a site foreman}

\label{lesson:KA-C306-mestri}



\begin{quote}
Before the new one: say the Kannada for a cobbler, then the Kannada for a potter.
\end{quote}

\subsection*{You'll want to know: \kn{ಮೇಸ್ತ್ರಿ}}
\textbf{\kn{ಮೇಸ್ತ್ರಿ}} — \emph{mēstri} — \textquotedblleft{}a mason, a site foreman\textquotedblright{}.

\begin{grammarlens}[title={What you've built}]
One more word for this chapter.
\end{grammarlens}

\subsection*{Your turn}
\begin{itemize}
\raggedright
  \item \emph{Say it:} \emph{mēstri}
  \item \emph{Say it:} \emph{mēstri}, once more
  \item \emph{Say it:} say \emph{mēstri}
  \item \emph{Recall:} say the Kannada for a cobbler, then the Kannada for a potter, then say \emph{mēstri} again
\end{itemize}

\begin{tcolorbox}[breakable,colback=teal!4,colframe=teal!35!black,title={Before you move on}]
What does \kn{ಮೇಸ್ತ್ರಿ} mean? (\textquotedblleft{}A mason, a site foreman\textquotedblright{}.) Say it once more.
\end{tcolorbox}
\section[\texorpdfstring{\kn{ಪ್ಲಂಬರ್} (plambar)}{ಪ್ಲಂಬರ್ (plambar)}]{\kn{ಪ್ಲಂಬರ್} (plambar) — a plumber}

\label{lesson:KA-C306-plambar}



\begin{quote}
Before the new one: say the Kannada for a potter, then the Kannada for a mason.
\end{quote}

\subsection*{You'll want to know: \kn{ಪ್ಲಂಬರ್}}
\textbf{\kn{ಪ್ಲಂಬರ್}} — \emph{plambar} — \textquotedblleft{}a plumber\textquotedblright{}.

\begin{grammarlens}[title={What you've built}]
One more word for this chapter.
\end{grammarlens}

\subsection*{Your turn}
\begin{itemize}
\raggedright
  \item \emph{Say it:} \emph{plambar}
  \item \emph{Say it:} \emph{plambar}, once more
  \item \emph{Say it:} say \emph{plambar}
  \item \emph{Recall:} say the Kannada for a potter, then the Kannada for a mason, then say \emph{plambar} again
\end{itemize}

\begin{tcolorbox}[breakable,colback=teal!4,colframe=teal!35!black,title={Before you move on}]
What does \kn{ಪ್ಲಂಬರ್} mean? (\textquotedblleft{}A plumber\textquotedblright{}.) Say it once more.
\end{tcolorbox}
\section[\texorpdfstring{\kn{ತೋಟಗಾರ} (tōṭagāra)}{ತೋಟಗಾರ (tōṭagāra)}]{\kn{ತೋಟಗಾರ} (tōṭagāra) — a gardener}

\label{lesson:KA-C306-totagara}



\begin{quote}
Before the new one: say the Kannada for a mason, then the Kannada for a plumber.
\end{quote}

\subsection*{You'll want to know: \kn{ತೋಟಗಾರ}}
\textbf{\kn{ತೋಟಗಾರ}} — \emph{tōṭagāra} — \textquotedblleft{}a gardener\textquotedblright{}.

\begin{grammarlens}[title={What you've built}]
One more word for this chapter.
\end{grammarlens}

\subsection*{Your turn}
\begin{itemize}
\raggedright
  \item \emph{Say it:} \emph{tōṭagāra}
  \item \emph{Say it:} \emph{tōṭagāra}, once more
  \item \emph{Say it:} say \emph{tōṭagāra}
  \item \emph{Recall:} say the Kannada for a mason, then the Kannada for a plumber, then say \emph{tōṭagāra} again
\end{itemize}

\begin{tcolorbox}[breakable,colback=teal!4,colframe=teal!35!black,title={Before you move on}]
What does \kn{ತೋಟಗಾರ} mean? (\textquotedblleft{}A gardener\textquotedblright{}.) Say it once more.
\end{tcolorbox}
